\section{\label{sec:method}Method}

\subsection{Problem formulation}

Given a device permittivity map $\varepsilon(x,y)$, a source excitation mask $m_\text{src}(x,y)$, and a free-space
wavelength $\lambda$, we seek the complex electric field $E_z(x,y)\in\mathbb{C}$ that solves the Helmholtz equation
[Eq.~\eqref{eq:helmholtz}]. We frame this as a conditional generation problem: learn a model that maps Gaussian noise
to the field solution, conditioned on $(\varepsilon, m_\text{src}, \lambda)$.

The model input is a tensor $\mathbf{x}\in\mathbb{R}^{4\times H\times W}$ whose channels are the real and imaginary
parts of $E_z$ (interpolated between noise and target during training), the permittivity map, and the source mask. The
wavelength enters as a scalar conditioning variable via the timestep embedding. The model outputs a two-channel
velocity field $\hat{u}_\theta\in\mathbb{R}^{2\times H\times W}$ representing the real and imaginary components of the
predicted flow-matching velocity.

\subsection{Complex-valued U-Net architecture}

We use a complex-valued U-Net that operates natively on split real/imaginary representations, preserving phase
relationships throughout the network. The architecture follows the standard encoder--bottleneck--decoder layout with
skip connections.

\paragraph{Complex convolution.}
Each convolutional layer implements full complex multiplication using two real weight
tensors $W_r$ and $W_i$~\cite{Trabelsi2018DeepComplex}:
\begin{align}
    y_r &= W_r * x_r - W_i * x_i, \nonumber\\
    y_i &= W_r * x_i + W_i * x_r,
    \label{eq:complex_conv}
\end{align}
where $*$ denotes spatial convolution. Complex biases are added independently to the real and imaginary outputs.
Weights are initialized with complex Glorot normal initialization: $\mathrm{Var}[W_r] = \mathrm{Var}[W_i] =
1/(f_\text{in} + f_\text{out})$, so that $\mathrm{Var}[W_r + iW_i] = 2/(f_\text{in} + f_\text{out})$.

\paragraph{ModReLU activation.}
We use the modulus ReLU~\cite{Arjovsky2016UnitaryRNN}, which gates the magnitude while
preserving phase:
\begin{equation}
    \mathrm{ModReLU}(z) = \mathrm{ReLU}(|z| + b) \cdot \frac{z}{|z|},
    \label{eq:modrelu}
\end{equation}
where $b$ is a learnable bias per channel, initialized to $-0.1$. This allows the network to suppress low-amplitude
features without distorting the phase of signals that pass through.

\paragraph{Complex group normalization.}
Real and imaginary parts are normalized jointly: both are centered by their
respective group means, then divided by the shared complex standard deviation $\sigma =
\sqrt{\mathrm{Var}[x_r]+\mathrm{Var}[x_i]+\epsilon}$. Learnable affine parameters $(\gamma_r, \gamma_i, \beta_r,
\beta_i)$ apply a complex scale-and-shift after normalization:
\begin{align}
    y_r &= \gamma_r \hat{x}_r - \gamma_i \hat{x}_i + \beta_r, \nonumber\\
    y_i &= \gamma_r \hat{x}_i + \gamma_i \hat{x}_r + \beta_i.
    \label{eq:complex_groupnorm}
\end{align}

\paragraph{Complex self-attention.}
At selected resolution levels, we apply multi-head self-attention with complex-valued
queries, keys, and values. Attention logits are computed from the real part of the Hermitian inner product
$\mathrm{Re}(QK^\dagger)$, scaled by $\sqrt{d_h}$, then softmaxed and applied to complex values:
\begin{equation}
    \mathrm{Attn}(Q,K,V) = \mathrm{softmax}\!\left(\frac{\mathrm{Re}(QK^\dagger)}{\sqrt{d_h}}\right) V.
    \label{eq:complex_attn}
\end{equation}

\paragraph{Conditioning.}
The flow-matching timestep $t$ is encoded via sinusoidal positional embeddings and concatenated
with the normalized wavelength $\lambda$. This combined embedding passes through a two-layer MLP to produce a
conditioning vector that modulates each residual block via complex scale injection: the embedding is projected to a
pair of real-valued scale maps $(s_r, s_i)$ that multiply the hidden activations as a complex number, i.e.\
$h \leftarrow (s_r + is_i) \cdot h + h$.

\paragraph{Physics embedding.}
The Helmholtz residual $R = \nabla^2 E_z + k_0^2 \varepsilon E_z$ is computed from the
current noisy field using a five-point Laplacian stencil and concatenated as an additional complex input channel to the
U-Net stem. This residual is masked to exclude source and PML regions, normalized by its spatial magnitude, and
time-gated (ramped on for $t > 0.2$) so that the physics signal is suppressed when the field is still dominated by
noise at early flow-matching times.

\paragraph{Auxiliary inputs.}
Real-valued conditioning maps (permittivity and source mask) are projected into complex
space via a learned \textsc{RealToComplex} bridge (a pair of real convolutions producing real and imaginary parts) and
concatenated with the complex field channels before the stem convolution.

\subsection{Training objective}

The total loss combines four terms:
\begin{equation}
    \mathcal{L} = \mathcal{L}_\text{FM}
    + \lambda_R\,\mathcal{L}_\text{res}
    + \lambda_\phi\,\mathcal{L}_\text{phase}
    + \lambda_S\,\mathcal{L}_\text{S-param},
    \label{eq:total_loss}
\end{equation}
with weights $\lambda_R$, $\lambda_\phi$, $\lambda_S$ that are ramped on according to the curriculum described in
Sec.~\ref{sec:curriculum}.

\paragraph{Flow matching loss.}
Given source--target pairs $(x_0, x_1)$ where $x_0 \sim \mathcal{N}(0,I)$ and $x_1$ is the
ground-truth field, we sample time $t$ from a 50/50 mixture of $\mathcal{U}[0,1]$ and a logit-normal distribution
$t = \sigma(\mu + s\cdot z)$, $z \sim \mathcal{N}(0,1)$, which concentrates samples near $t = 0.5$ where the
denoising task is hardest. The interpolated state is $x_t = (1-t)x_0 + t\,x_1$ and the target velocity is
$u_t = x_1 - x_0$. The loss is the spatially weighted MSE between the predicted and target velocities:
\begin{equation}
    \mathcal{L}_\text{FM} = \frac{\sum_{p} w(p)\,\|\hat{u}_\theta(x_t,t)(p) - u_t(p)\|^2}{\sum_{p} w(p)},
    \label{eq:fm_loss_weighted}
\end{equation}
where the spatial weight $w(p)$ combines a PML-exclusion mask with an optional device-region upweight.

\paragraph{Helmholtz residual loss.}
From the predicted velocity we reconstruct the clean field estimate
$\hat{x}_1 = x_t + (1-t)\hat{u}_\theta$ and de-normalize to physical units. The Helmholtz residual is
\begin{equation}
    R(p) = \nabla^2 \hat{E}_z(p) + k_0^2\,\varepsilon(p)\,\hat{E}_z(p),
    \label{eq:residual}
\end{equation}
computed with a five-point finite-difference Laplacian. The loss is the masked, normalized mean-squared residual:
\begin{equation}
    \mathcal{L}_\text{res} = \frac{\sum_p\, w_R(p)\,|R(p)|^2}{\sum_p\, w_R(p)\,|k_0^2\varepsilon(p) E_z^\text{true}(p)|^2},
    \label{eq:residual_loss}
\end{equation}
where the weight mask $w_R$ is the product of three binary masks: (i) a PML-exclusion mask that zeros out absorbing
boundary cells, (ii) a source-exclusion mask (the Helmholtz equation is only valid in source-free regions), and (iii)
an interface mask that excludes pixels where $|\nabla\varepsilon|$ exceeds a threshold, since the five-point Laplacian
stencil produces spurious residuals at sharp dielectric boundaries (Si/SiO$_2$). The denominator normalizes by the
driving-term scale $|k_0^2\varepsilon E_z|^2$, making the loss approximately 0 when physics is satisfied and 1 when
the residual equals the driving term, independent of grid spacing and field amplitude. When time-gating is enabled,
samples with $t$ below a threshold are excluded from the residual computation, since the reconstructed field at low $t$
is too noisy to yield informative physics gradients.

\paragraph{Phase loss.}
We penalize phase disagreement between the predicted and ground-truth fields using a
wrap-safe cosine distance:
\begin{equation}
    \mathcal{L}_\text{phase} = 1 - \mathrm{Re}\!\left(\frac{\hat{E}_z}{|\hat{E}_z|} \cdot
    \frac{\bar{E}_z^\text{true}}{|E_z^\text{true}|}\right),
    \label{eq:phase_loss}
\end{equation}
weighted by $|E_z^\text{true}|^2$ and gated by an amplitude threshold so that phase is only supervised in regions with
appreciable field strength. Before computing the phase loss, predicted and true fields are globally phase-aligned via a
weighted inner product to remove the gauge freedom inherent in complex field solutions.

\paragraph{S-parameter loss.}
S-parameters are extracted from the predicted field by projecting onto the fundamental
waveguide eigenmode at each port cross-section. At port $j$, the mode amplitude is
\begin{equation}
    a_j = \int \hat{E}_z(r_\perp)\,\bar{e}_j(r_\perp)\,dr_\perp,
    \label{eq:overlap}
\end{equation}
where $e_j$ is the pre-computed mode profile obtained by solving the 1D transverse Helmholtz eigenvalue problem at the
port location and operating wavelength, and the integral runs over the port cross-section. The S-parameter is then
$S_{ji} = a_j / a_i$ where $i$ is the excited input port. The loss penalizes both magnitude and phase errors relative
to FDTD-extracted ground truth:
\begin{equation}
    \mathcal{L}_\text{S-param} = \frac{1}{|\mathcal{P}|}\sum_{j\in\mathcal{P}} g_j\!\left[
        w_\text{mag}\,(|S_{j}^\text{pred}| - |S_{j}^\text{true}|)^2
        + w_\phi\left(1 - \mathrm{Re}\!\left(\frac{S_j^\text{pred}}{|S_j^\text{pred}|}
          \cdot \frac{\bar{S}_j^\text{true}}{|S_j^\text{true}|}\right)\right)
    \right],
    \label{eq:sparam_loss}
\end{equation}
where $g_j = \mathbf{1}[|S_j^\text{true}| > \tau]$ is an amplitude gate that suppresses supervision on near-zero
transmissions, the sum runs over valid output ports $\mathcal{P}$ (excluding the excited input port, whose total-field
projection conflates incident and scattered waves), and the denominator counts only gated ports for a proper masked
mean.

\subsection{\label{sec:curriculum}Multi-loss curriculum}

The four loss terms are introduced in three phases to avoid gradient conflicts between competing objectives:

\begin{itemize}
    \item \textbf{Phase~A} (epochs $0$--$E_A$): Flow matching loss only. The model learns the basic transport from
          noise to fields.
    \item \textbf{Phase~B} (epochs $E_A$--$E_B$): Helmholtz residual loss is activated with a linear warmup over
          $W_R$ epochs. The model begins enforcing wave-equation compliance on its predictions.
    \item \textbf{Phase~C} (epochs $E_B$--end): Phase loss and S-parameter loss are activated, each with independent
          linear warmups. The model refines phase accuracy and port-to-port coupling.
\end{itemize}

Within each phase, the loss weights ramp linearly from 0 to their target values over a configurable warmup period.
Optionally, conflict-free gradient surgery~\cite{ConFIG2025} is applied after Phase~B to project conflicting
per-loss gradients into a common descent direction, preventing destructive interference between the velocity-matching
and physics-compliance objectives.

\subsection{Inference}

At inference time, we generate field predictions by integrating the learned velocity field from $t = 0$ (Gaussian
noise) to $t = 1$ using a second-order Heun (improved Euler) integrator:
\begin{align}
    \tilde{x}_{k+1} &= x_k + \Delta t_k\, \hat{u}_\theta(x_k, t_k), \nonumber\\
    x_{k+1} &= x_k + \tfrac{1}{2}\Delta t_k\!\left[\hat{u}_\theta(x_k, t_k)
               + \hat{u}_\theta(\tilde{x}_{k+1}, t_{k+1})\right].
    \label{eq:heun}
\end{align}
The time grid uses a quadratic schedule $t_k = (k/N)^2$ that allocates more steps near $t = 0$ where the field
is noisiest and the velocity field changes most rapidly. S-parameters are extracted from the final field via the same
modal overlap integral used during training, requiring no additional learned components.
